% SPDX-FileCopyrightText: 2025 IObundle
%
% SPDX-License-Identifier: MIT

\begin{itemize}
  \itemsep-0.5em
  \item \textbf{Dataflow computation}: an accelerator is described as data
    moving across units that operate on it as it passes through; units are
    sources, compute units, or sinks, depending on whether they only
    produce, transform, or consume data.
  \item \textbf{Composability}: complex units are built by instantiating
    and interconnecting simpler, previously defined units. A unit can never
    instantiate itself, directly or indirectly, since recursive
    instantiation is not permitted in a dataflow design.
  \item \textbf{Bounded-latency units}: Versat only supports units that
    take a fixed amount of time to process their input, which is what
    allows it to reason about data validity and timing across the whole
    accelerator automatically.
  \item \textbf{Hardware/software co-generation}: hardware and software are
    never designed separately; both are derived from the same
    specification, so the generated driver is always in sync with the
    generated accelerator.
  \item \textbf{Extensibility}: designers can supply their own custom basic
    units, written directly in Verilog, as long as they implement the
    interfaces Versat requires to integrate them with the rest of the
    generated accelerator.
\end{itemize}
